% GENERATED FILE. Edit canonical lessons, then run npm run generate:books.
% canonical-source-hash: fnv1a64:e10b4851b662d82c
% canonical-lessons: FA-C350-padeshah, FA-C350-malake, FA-C350-kakh, FA-C350-qale, FA-C350-taj

\chapter{Kings and castles}
\label{ch:fa-kings-and-castles}
\hlchaptermodality{350}

\begin{chapteropening}
\textbf{By the end of this chapter:} \emph{I can talk about a king, a queen, a palace, a castle and a crown.}

Complete the last lesson of chapter 350: I can talk about a king, a queen, a palace, a castle and a crown.
\end{chapteropening}

\section[\texorpdfstring{\fa{پادشاه} (pâdeshâh)}{پادشاه (pâdeshâh)}]{\fa{پادشاه} (pâdeshâh) — a king, a monarch}

\label{lesson:FA-C350-padeshah}



\begin{quote}
Before the new one: say the Persian for to kill, then the Persian for to be taken prisoner.
\end{quote}

\subsection*{You'll want to know: \fa{پادشاه}}
\textbf{\fa{پادشاه}} (\emph{pâdeshâh}) — \textquotedblleft{}a king, a monarch\textquotedblright{}.

\begin{quote}
\fa{پادشاه} \fa{سه} \fa{پسر} \fa{داشت}.
\end{quote}

\begin{quote}
\emph{Pâdeshâh se pesar dâsht.}
\end{quote}

\begin{quote}
The king had three sons.
\end{quote}

Also \emph{shâh}. \emph{Pâdeshâhi} \textquotedblleft{}a kingdom, a reign\textquotedblright{}.

\begin{cousinweb}
\emph{pâdeshâh} is a native Persian word that EtymDB glosses \textquotedblleft{}king, monarch\textquotedblright{} without a parent; it is usually traced to Middle Persian \emph{pādixšāh}, \textquotedblleft{}ruler\textquotedblright{}.
\end{cousinweb}

\begin{grammarlens}[title={What you've built}]
One more word for telling, explaining and describing at B1.
\end{grammarlens}

\subsection*{Your turn}
\begin{itemize}
\raggedright
  \item \emph{Say it:} \emph{pâdeshâh}
  \item \emph{Say it:} \emph{pâdeshâh}, once more
  \item \emph{Say it:} a sentence of your own with \emph{pâdeshâh}
  \item \emph{Recall:} say the Persian for to kill, then the Persian for to be taken prisoner, then say \emph{pâdeshâh} again
\end{itemize}

\begin{tcolorbox}[breakable,colback=teal!4,colframe=teal!35!black,title={Before you move on}]
What does \fa{پادشاه} mean? (\textquotedblleft{}A king, a monarch\textquotedblright{}.) Say it once more.
\end{tcolorbox}
\section[\texorpdfstring{\fa{ملکه} (malake)}{ملکه (malake)}]{\fa{ملکه} (malake) — a queen}

\label{lesson:FA-C350-malake}



\begin{quote}
Before the new one: say the Persian for to be taken prisoner, then the Persian for a king.
\end{quote}

\subsection*{You'll want to know: \fa{ملکه}}
\textbf{\fa{ملکه}} (\emph{malake}) — \textquotedblleft{}a queen\textquotedblright{}.

\begin{quote}
\fa{ملکه} \fa{یک} \fa{دختر} \fa{زیبا} \fa{داشت}.
\end{quote}

\begin{quote}
\emph{Malake yek dokhtar-e zibâ dâsht.}
\end{quote}

\begin{quote}
The queen had a beautiful daughter.
\end{quote}

From the Arabic feminine of \emph{malek} \textquotedblleft{}king\textquotedblright{}.

\begin{cousinweb}
\emph{malake} is borrowed from Arabic \emph{malika}, \textquotedblleft{}queen\textquotedblright{} (EtymDB).
\end{cousinweb}

\begin{grammarlens}[title={What you've built}]
One more word for telling, explaining and describing at B1.
\end{grammarlens}

\subsection*{Your turn}
\begin{itemize}
\raggedright
  \item \emph{Say it:} \emph{malake}
  \item \emph{Say it:} \emph{malake}, once more
  \item \emph{Say it:} a sentence of your own with \emph{malake}
  \item \emph{Recall:} say the Persian for to be taken prisoner, then the Persian for a king, then say \emph{malake} again
\end{itemize}

\begin{tcolorbox}[breakable,colback=teal!4,colframe=teal!35!black,title={Before you move on}]
What does \fa{ملکه} mean? (\textquotedblleft{}A queen\textquotedblright{}.) Say it once more.
\end{tcolorbox}
\section[\texorpdfstring{\fa{کاخ} (kâkh)}{کاخ (kâkh)}]{\fa{کاخ} (kâkh) — a palace}

\label{lesson:FA-C350-kakh}



\begin{quote}
Before the new one: say the Persian for a king, then the Persian for a queen.
\end{quote}

\subsection*{You'll want to know: \fa{کاخ}}
\textbf{\fa{کاخ}} (\emph{kâkh}) — \textquotedblleft{}a palace\textquotedblright{}.

\begin{quote}
\fa{کاخ} \fa{پادشاه} \fa{روی} \fa{تپه} \fa{بود}.
\end{quote}

\begin{quote}
\emph{Kâkh-e pâdeshâh ru-ye tappe bud.}
\end{quote}

\begin{quote}
The king's palace stood on a hill.
\end{quote}

Plural \emph{kâkh-hâ}; many old palaces are museums now.

\begin{cousinweb}
\emph{kâkh} is a native Persian word for a palace. EtymDB lists no entry, and its deeper origin is not settled.
\end{cousinweb}

\begin{grammarlens}[title={What you've built}]
One more word for telling, explaining and describing at B1.
\end{grammarlens}

\subsection*{Your turn}
\begin{itemize}
\raggedright
  \item \emph{Say it:} \emph{kâkh}
  \item \emph{Say it:} \emph{kâkh}, once more
  \item \emph{Say it:} a sentence of your own with \emph{kâkh}
  \item \emph{Recall:} say the Persian for a king, then the Persian for a queen, then say \emph{kâkh} again
\end{itemize}

\begin{tcolorbox}[breakable,colback=teal!4,colframe=teal!35!black,title={Before you move on}]
What does \fa{کاخ} mean? (\textquotedblleft{}A palace\textquotedblright{}.) Say it once more.
\end{tcolorbox}
\section[\texorpdfstring{\fa{قلعه} (qal'e)}{قلعه (qal'e)}]{\fa{قلعه} (qal'e) — a castle, a fortress}

\label{lesson:FA-C350-qale}



\begin{quote}
Before the new one: say the Persian for a queen, then the Persian for a palace.
\end{quote}

\subsection*{You'll want to know: \fa{قلعه}}
\textbf{\fa{قلعه}} (\emph{qal'e}) — \textquotedblleft{}a castle, a fortress\textquotedblright{}.

\begin{quote}
\fa{دشمن} \fa{نتوانست} \fa{وارد} \fa{قلعه} \fa{شود}.
\end{quote}

\begin{quote}
\emph{Doshman natavânest vâred-e qal'e shavad.}
\end{quote}

\begin{quote}
The enemy could not get into the castle.
\end{quote}

Written with \fa{ع}. Plural \emph{qal'e-hâ}.

\begin{cousinweb}
\emph{qal'e} is borrowed from Arabic \emph{qal'a}, \textquotedblleft{}fort, fortress\textquotedblright{} (EtymDB), and EtymDB in turn connects the Arabic word with Persian \emph{kalât} \textquotedblleft{}fortress\textquotedblright{}.
\end{cousinweb}

\begin{grammarlens}[title={What you've built}]
One more word for telling, explaining and describing at B1.
\end{grammarlens}

\subsection*{Your turn}
\begin{itemize}
\raggedright
  \item \emph{Say it:} \emph{qal'e}
  \item \emph{Say it:} \emph{qal'e}, once more
  \item \emph{Say it:} a sentence of your own with \emph{qal'e}
  \item \emph{Recall:} say the Persian for a queen, then the Persian for a palace, then say \emph{qal'e} again
\end{itemize}

\begin{tcolorbox}[breakable,colback=teal!4,colframe=teal!35!black,title={Before you move on}]
What does \fa{قلعه} mean? (\textquotedblleft{}A castle, a fortress\textquotedblright{}.) Say it once more.
\end{tcolorbox}
\section[\texorpdfstring{\fa{تاج} (tâj)}{تاج (tâj)}]{\fa{تاج} (tâj) — a crown}

\label{lesson:FA-C350-taj}



\begin{quote}
Before the new one: say the Persian for a palace, then the Persian for a castle.
\end{quote}

\subsection*{You'll want to know: \fa{تاج}}
\textbf{\fa{تاج}} (\emph{tâj}) — \textquotedblleft{}a crown\textquotedblright{}.

\begin{quote}
\fa{پادشاه} \fa{یک} \fa{تاج} \fa{طلا} \fa{بر} \fa{سر} \fa{داشت}.
\end{quote}

\begin{quote}
\emph{Pâdeshâh yek tâj-e talâ bar sar dâsht.}
\end{quote}

\begin{quote}
The king wore a gold crown on his head.
\end{quote}

\emph{Tâj-o takht} \textquotedblleft{}crown and throne\textquotedblright{}, royal power.

\begin{cousinweb}
\emph{tâj} is an Iranian word: EtymDB traces it to Parthian \emph{tāg}, \textquotedblleft{}crown\textquotedblright{}, and lists Arabic \emph{tāj} beside it.
\end{cousinweb}

\begin{grammarlens}[title={What you've built}]
One more word for telling, explaining and describing at B1.
\end{grammarlens}

\subsection*{Your turn}
\begin{itemize}
\raggedright
  \item \emph{Say it:} \emph{tâj}
  \item \emph{Say it:} \emph{tâj}, once more
  \item \emph{Say it:} a sentence of your own with \emph{tâj}
  \item \emph{Recall:} say the Persian for a palace, then the Persian for a castle, then say \emph{tâj} again
\end{itemize}

\begin{tcolorbox}[breakable,colback=teal!4,colframe=teal!35!black,title={Before you move on}]
What does \fa{تاج} mean? (\textquotedblleft{}A crown\textquotedblright{}.) Say it once more.
\end{tcolorbox}
